\documentclass[9pt,a4paper]{article}

\usepackage[utf8]{inputenc}
\usepackage[T1]{fontenc}
\usepackage{mathpazo}
\usepackage[scaled=0.92]{helvet}
\usepackage[margin=0.5in,top=0.4in,bottom=0.4in]{geometry}
\usepackage{amsmath,amssymb}
\usepackage{booktabs}
\usepackage{tabularx}
\usepackage{multicol}
\usepackage{enumitem}
\usepackage{setspace}

\setstretch{1.0}
\setlength{\parindent}{0pt}
\setlength{\parskip}{0.2em}
\setlength{\columnsep}{0.4cm}
\setlength{\columnseprule}{0.3pt}
\setlength{\emergencystretch}{2em}

\setlist[itemize]{leftmargin=1.2em,noitemsep,topsep=1pt,parsep=0pt,itemsep=1pt}
\setlist[enumerate]{leftmargin=1.2em,noitemsep,topsep=1pt,parsep=0pt,itemsep=1pt}
\renewcommand{\arraystretch}{1.15}

\newcommand{\mysect}[1]{%
  \par\vspace{5pt}%
  \noindent{\small\bfseries\sffamily #1}\par\vspace{-2pt}%
  \noindent\rule{\linewidth}{0.4pt}\par\vspace{2pt}%
  \noindent\ignorespaces%
}

\pagestyle{empty}

\begin{document}

\begin{center}
	{\Large\bfseries\sffamily Cheat Sheet: Relasi}\\[2pt]
	{\small Matematika Diskrit \textbullet\ TRPL Politeknik Negeri Batam \textbullet\ Pertemuan 4}
\end{center}
\vspace{-6pt}

\begin{multicols}{2}

	\mysect{Cartesian Product}
	$A \times B = \{(a,b) \mid a \in A,\ b \in B\}$.\\
	Kardinalitas: $|A \times B| = |A| \cdot |B|$.\\
	Contoh: $A=\{a,b\}, B=\{1,2,3\}$.\\
	$A \times B = \{(a,1),(a,2),(a,3),(b,1),(b,2),(b,3)\}$.\\
	Catatan: $A \times B \neq B \times A$.

	\mysect{Definisi Relasi}
	$R \subseteq A \times B$.\\
	Domain $= A$, Kodomain $= B$.\\
	Range $= \mathrm{ran}(R) = \{b \in B \mid (a,b) \in R\}$.\\
	Relasi khusus: $\emptyset$ (kosong), $A \times B$ (universal), $I_A$ (identitas).

	\mysect{Representasi}
	Misal $A=\{a,b\}, B=\{1,2,3\}, R=\{(a,1),(b,3)\}$.
	\begin{enumerate}
		\item \textbf{Diagram Panah:} $a \to 1$, $b \to 3$.
		\item \textbf{Tabel:} Pasangan $(a,1)$ dan $(b,3)$.
		\item \textbf{Matriks Biner:} $m_{ij}=1$ jika terhubung, $0$ jika tidak.
		      \[ M_R = \begin{pmatrix} 1 & 0 & 0 \\ 0 & 0 & 1 \end{pmatrix} \]
	\end{enumerate}

	\mysect{Sifat Relasi pada Himpunan $A$}
	{\small
		\begin{tabular}{@{}l l l@{}}
			\toprule
			Sifat        & Syarat                                & Cek $M_R$        \\
			\midrule
			Refleksif    & $(a,a) \in R, \forall a$              & $m_{ii}=1$       \\
			Irefleksif   & $(a,a) \notin R, \forall a$           & $m_{ii}=0$       \\
			Simetris     & $(a,b) \in R \Rightarrow (b,a) \in R$ & $M = M^T$        \\
			Antisimetris & $(a,b),(b,a) \in R \Rightarrow a=b$   & $m_{ij}m_{ji}=0$ \\
			Transitif    & $(a,b),(b,c) \in R \Rightarrow (a,c)$ & $M^2 \le M$      \\
			\bottomrule
		\end{tabular}
	}
	\vspace{2pt}\\
	\textbf{Penting:} Tidak simetris $\neq$ antisimetris.

	\mysect{Invers Relasi}
	$R^{-1} = \{(b,a) \mid (a,b) \in R\}$.\\
	Matriks: $M_{R^{-1}} = M_R^T$ (\textbf{bukan} $M_R^{-1}$!).\\
	Contoh: $R=\{(a,1),(b,3)\} \Rightarrow R^{-1}=\{(1,a),(3,b)\}$.\\
	Sifat: $(R^{-1})^{-1} = R$.

	\mysect{Kombinasi Relasi}
	Untuk $R_1, R_2$ dari $A$ ke $B$:
	\begin{itemize}
		\item $R_1 \cap R_2$: Irisan
		\item $R_1 \cup R_2$: Gabungan
		\item $R_1 - R_2$: Selisih
		\item $R_1 \oplus R_2$: Beda Simetris
	\end{itemize}

	\mysect{Komposisi Relasi}
	$S \circ R = \{(a,c) \mid \exists b, (a,b) \in R \land (b,c) \in S\}$.\\
	Dibaca \textbf{kanan ke kiri}: kerjakan $R$ dulu, lalu $S$.\\
	Contoh: $R=\{(1,2),(2,3)\}, S=\{(2,a),(3,b)\}$.\\
	$S \circ R = \{(1,a), (2,b)\}$.\\
	Sifat: Asosiatif $T \circ (S \circ R) = (T \circ S) \circ R$.

	\mysect{Studi Kasus: Sistem Akademik}
	$R_1$ (Mengambil) $= \{(M1,MTK), (M2,RPL), (M3,MTK), (M3,RPL)\}$.\\
	$R_2$ (Mengampu) $= \{(MTK,D1), (RPL,D2)\}$.
	\begin{itemize}
		\item \textbf{Pengambil MTK:} Gunakan $R_1^{-1}$ pada MTK $\Rightarrow \{M1, M3\}$.
		\item \textbf{Dosen M3:} Gunakan $R_2 \circ R_1$. Jalur: $M3 \to MTK \to D1$, dan $M3 \to RPL \to D2$. Hasil $\Rightarrow \{D1, D2\}$.
	\end{itemize}

	\mysect{Fungsi}
	Relasi $f: A \to B$ dengan dua syarat mutlak:
	\begin{enumerate}
		\item Setiap $a \in A$ punya pasangan (total).
		\item Pasangan tunggal (well-defined).
	\end{enumerate}

	\mysect{Kesalahan Umum}
	{\small
		\begin{tabular}{@{}l l@{}}
			\toprule
			Salah                           & Benar                        \\
			\midrule
			$A \times B = B \times A$       & $A \times B \neq B \times A$ \\
			Invers relasi $= M^{-1}$        & Invers relasi $= M^T$        \\
			Tidak simetris $=$ antisimetris & Dua hal berbeda              \\
			Range $=$ kodomain              & Range $\subseteq$ kodomain   \\
			\bottomrule
		\end{tabular}
	}

	\mysect{Lembar Rujukan Cepat}
	{\small
		\begin{tabular}{@{}l l@{}}
			\toprule
			Butuh                     & Rumus / Tipe              \\
			\midrule
			Kardinalitas $A \times B$ & $|A| \cdot |B|$           \\
			Cek Refleksif             & Diagonal $M_R = 1$        \\
			Cek Simetris              & $M_R = M_R^T$             \\
			Cek Transitif             & $M_R^2 \le M_R$           \\
			Invers Relasi             & $R^{-1}$, matriks $M_R^T$ \\
			Komposisi                 & $S \circ R$ (baca kanan)  \\
			Relasi Universal          & $R = A \times B$          \\
			Relasi Identitas          & $I_A = \{(a,a)\}$         \\
			\bottomrule
		\end{tabular}
	}

\end{multicols}

\end{document}
